\section{Control Laws in Robotics}

\subsection{Proportional-Derivative (PD)}
\label{subsec:pd_control}

Proportional-Derivative (PD) control is the cornerstone of robotic motion control. In Robotics 2,
 we move beyond simple linear systems to consider the full non-linear dynamics of an $n$-DOF manipulator, typically expressed as:
\begin{equation}
    B(\bm{q})\ddot{\bm{q}} + \bm{C}(\bm{q}, \dot{\bm{q}})\dot{\bm{q}} + \bm{g}(\bm{q}) = \bm{\tau}
\end{equation}
where $B(\bm{q})$ is the inertia matrix, $\bm{C}$ represents Coriolis/centrifugal terms, and $\bm{g}$ is the gravity vector. 
PD control is used to compute the torque $\bm{\tau}$ required to achieve a specific task.
It calculates control output based on both the current error and the rate of change of that error, making it effective for systems
requiring fast response and minimal overshoot.
It functions as a tracker for a system to a reference point.

\subsubsection{Regulation and the Mechanical Analogy}
For a regulation task (reaching a constant set-point $\bm{q}_d$), the PD control law is:
\begin{equation}
    \bm{\tau} = \bm{K}_P (\bm{q}_d - \bm{q}) - \bm{K}_D \dot{\bm{q}}
\end{equation}
This can be physically interpreted as an \textbf{Impedance Control} strategy. 

The controller shapes the robot's behavior so that it mimics a passive mechanical system where the joints are connected to $\bm{q}_d$ via virtual springs ($\bm{K}_P$)
and virtual dampers ($\bm{K}_D$). This analogy is crucial because it ensures that, in the absence of gravity, the system is inherently stable due to its passive nature.

\begin{framed}
\noindent \textit{Notes: Impedance control is a robotics strategy that regulates 
the dynamic interaction between a robot and its environment by defining a desired relationship between applied 
force and resulting motion (mass-spring-damper system). The \textbf{Core Principle} is that the robot's end-effector is controlled to follow a 
specified impedance behavior, governed by the following dynamic equation:
\begin{equation}
    \bm{F} = \bm{M}\bm{a} + \bm{C}\bm{v} + \bm{K}\bm{x} + \bm{f} + \bm{s}
\end{equation}
where $\bm{F}$ is the interaction force, $\bm{x}$ is the position, $\bm{v}$ the velocity, and $\bm{a}$ the acceleration. 
The parameters $\bm{K}$, $\bm{C}$, and $\bm{M}$ represent the stiffness, damping, and inertial matrices respectively, 
while $\bm{f}$ accounts for friction and $\bm{s}$ for static forces.}
\end{framed}

\subsubsection{Stability Analysis via Lyapunov}
A fundamental result in Robotics 2 is the stability proof for PD control. 
By choosing a Lyapunov candidate function based on the total energy of the system (kinetic + potential error):
\begin{equation}
    V(\bm{e}, \dot{\bm{q}}) = \frac{1}{2}\dot{\bm{q}}^T B(\bm{q})\dot{\bm{q}} + \frac{1}{2}\bm{e}^T \bm{K}_P \bm{e}
\end{equation}
it can be shown that $\dot{V} = -\dot{\bm{q}}^T \bm{K}_D \dot{\bm{q}} \leq 0$. 
Using LaSalle's Invariance Principle, we conclude that the system is \textbf{Asymptotically Stable}, 
meaning the robot will eventually stop at the desired configuration $\bm{q}_d$, provided that gravity is compensated.



[Image of a mass-spring-damper system model]


\subsubsection{PD with Gravity Compensation}
The primary drawback of a pure PD controller is the presence of a steady-state error $\bm{e}_{\infty}$ 
caused by the gravity vector $\bm{g}(\bm{q})$. To achieve zero error without using an Integral (I) term—which can introduce instability—we employ 
\textbf{PD with Gravity Compensation}:
\begin{equation}
    \bm{\tau} = \bm{K}_P \bm{e} - \bm{K}_D \dot{\bm{q}} + \bm{g}(\bm{q})
\end{equation}
This model-based approach "cancels" the gravitational pull at every configuration, effectively making the robot behave as if it were in a zero-G environment.

\subsubsection{Extension to Task Space Control}
In many applications, the control objective is defined in the operational space (Cartesian coordinates $\bm{x}$) rather than joint space. 
The PD law is then mapped through the \textbf{Analytic Jacobian} $\bm{J}(\bm{q})$:
\begin{equation}
    \bm{\tau} = \bm{J}^T(\bm{q}) \left[ \bm{K}_P (\bm{x}_d - \bm{x}) - \bm{K}_D \dot{\bm{x}} \right] + \bm{g}(\bm{q})
\end{equation}
This allows the controller to define the stiffness and damping directly at the end-effector level, which is essential for interaction tasks and human-robot collaboration.


\subsection{Proportional-Integral-Derivative (PID)}
\label{subsec:pid_control}

The PD control law previously discussed exhibits a fundamental limitation: it cannot eliminate the \textbf{steady-state error} 
($e_{\infty}$) when the system is subject to persistent, unmodeled forces. In robotics, these typically include gravity, joint friction,
or sensor calibration offsets. Since a PD controller requires a non-zero error to generate a static compensation torque, the system will settle at 
an equilibrium point where the proportional action exactly balances the external disturbance, but $e \neq 0$.

To address this, an \textbf{Integral (I)} term is added, which considers the cumulative sum of past errors. The resulting PID control law in the time domain is:
\begin{equation}
    u(t) = K_P e(t) + K_I \int_{0}^{t} e(\tau) d\tau + K_D \dot{e}(t)
\end{equation}
where $K_I$ is the integral gain. In the Laplace domain, the transfer function of the controller is expressed as:
\begin{equation}
    K(s) = K_P + \frac{K_I}{s} + K_D s
\end{equation}

\subsubsection{Role of the Integral Term}
The integral action provides a "memory" for the controller. As long as a residual error exists, the integral term will continue to grow, 
increasing the control output until the error is driven to zero. This effectively provides \textbf{automatic compensation} for constant disturbances 
without requiring an explicit model of the force (e.g., without needing to know the exact value of the gravitational vector $g(q)$).

\begin{framed}
\noindent \textit{Notes: While the integral term eliminates steady-state error, it can introduce instability. 
In robotics, excessive $K_I$ leads to \textbf{overshoot} and oscillations. Furthermore, practical implementations must account 
for \textbf{Integrator Wind-up}, a phenomenon where the integral term grows excessively when the actuators saturate, 
leading to large overshoots when the error finally changes sign.}
\end{framed}

\subsubsection{Summary of Control Modes}
The synergy of the three terms allows for comprehensive control:
\begin{itemize}
    \item \textbf{Proportional (P):} Provides immediate corrective action based on the current error magnitude.
    \item \textbf{Integral (I):} Eliminates steady-state offsets by considering the error history.
    \item \textbf{Derivative (D):} Predicts future error trends to dampen the response and improve stability.
\end{itemize}